\chapter{Метод дифференцируемого рендеринга для реконструкции пространственно неоднородного площадного источника}

\section{Постановка задачи}

Рассматривается обратная задача реконструкции параметров пространственно неоднородного прямоугольного площадного источника освещения по HDR-изображениям известной сцены. Сцена образована тремя взаимно перпендикулярными плоскостями: горизонтальным полом и двумя вертикальными стенами, образующими угол. Все поверхности сцены имеют известный нейтральный серый диффузный материал.

Источник расположен над сценой в горизонтальной плоскости. Геометрические параметры источника считаются известными и не оптимизируются:
\begin{equation}
\mathcal{G}_L=
\left\{
\mathbf{c}_L,\,
 a,\,
 b,\,
 \mathbf{t}_1,\,
 \mathbf{t}_2,\,
 \mathbf{n}_L
\right\},
\end{equation}
где $\mathbf{c}_L \in \mathbb{R}^3$ --- положение центра источника; $a$ и $b$ --- длина и ширина прямоугольной светящейся области; $\mathbf{t}_1$ и $\mathbf{t}_2$ --- локальные ортогональные направления вдоль сторон источника; $\mathbf{n}_L=(0,0,-1)^\mathsf{T}$ --- нормаль источника, направленная вниз.

Неизвестными являются пространственно изменяющиеся параметры эмиссии, заданные в UV-пространстве источника пятканальной текстурой:
\begin{equation}
\mathbf{T}(u,v)=
\left[
\theta(u,v),\,
\eta(u,v),\,
c_R(u,v),\,
c_G(u,v),\,
c_B(u,v)
\right],
\qquad
(u,v)\in[0,1]^2.
\end{equation}

Здесь $\theta(u,v)$ --- параметр угловой ширины излучения; $\eta(u,v)\geq 0$ --- интенсивность излучения; $\mathbf{c}(u,v)=[c_R,c_G,c_B]^\mathsf{T}$ --- нормированный цвет локального элемента источника.

Цель реконструкции состоит в нахождении такой текстуры $\hat{\mathbf{T}}$, для которой HDR-изображения, синтезированные по известной геометрии сцены и восстановленному источнику, согласуются с набором наблюдаемых HDR-изображений:
\begin{equation}
\hat{\mathbf{T}}
=
\arg\min_{\mathbf{T}}
\mathcal{L}
\left(
\mathcal{R}(\mathcal{S},\mathcal{G}_L,\mathbf{T}),
\mathbf{H}^{\ast}
\right),
\end{equation}
где $\mathcal{S}$ --- известная геометрия и материалы сцены; $\mathcal{R}$ --- дифференцируемый рендерер; $\mathbf{H}^{\ast}$ --- наблюдаемые HDR-кадры в линейном цветовом пространстве.

\section{Параметризация источника}

\subsection{Дискретизация по площади}

Светящаяся поверхность разбивается на регулярную сетку размера $W_L\times H_L$. Каждая ячейка $i$ соответствует малому площадному элементу $\Delta A_i$ с центром $\mathbf{x}_i$:
\begin{equation}
\mathbf{x}_i
=
\mathbf{c}_L
+
\left(u_i-\frac{1}{2}\right)a\mathbf{t}_1
+
\left(v_i-\frac{1}{2}\right)b\mathbf{t}_2.
\end{equation}

Для каждого элемента оптимизируется набор параметров:
\begin{equation}
\boldsymbol{\pi}_i=
\left\{
\theta_i,\,
\eta_i,\,
\mathbf{c}_i
\right\}.
\end{equation}

Следовательно, число степеней свободы на уровне текстуры $H_L\times W_L$ равно
\begin{equation}
5H_LW_L.
\end{equation}

Такая дискретизация согласуется с конструкцией реальных осветителей, состоящих из набора отдельных светодиодов или LED-кластеров. В этой интерпретации интенсивность и цвет соседних ячеек описывают пространственное распределение вкладов отдельных светодиодов, а параметр $\theta_i$ приближённо моделирует локальную направленность каждого LED-элемента.

Если светодиодная матрица закрыта матовой рассеивающей крышкой, истинное распределение излучения по поверхности становится более гладким. В предлагаемой модели это учитывается не явным моделированием толщины, рассеяния и подповерхностного переноса света в диффузоре, а регуляризацией пространственных полей источника и реконструкцией на ограниченном разрешении. Таким образом, LED-панель с матовой крышкой рассматривается как эквивалентный площадной излучатель с гладкой эффективной текстурой параметров.

\subsection{Неотрицательная интенсивность}

Интенсивность параметризуется неограниченной переменной $z_i^\eta$:
\begin{equation}
\eta_i
=
\operatorname{softplus}(z_i^\eta)
=
\log\left(1+\exp(z_i^\eta)\right).
\end{equation}

Это обеспечивает физическое ограничение
\begin{equation}
\eta_i\geq 0
\end{equation}
без необходимости жёсткой проекции после каждого шага оптимизации. По сравнению с экспоненциальной параметризацией функция $\operatorname{softplus}$ уменьшает риск чрезмерного роста градиентов при реконструкции ярких компонентов источника.

\subsection{Нормированный цвет}

Цвет каждой ячейки задаётся вектором неограниченных параметров $\mathbf{z}_i^c \in \mathbb{R}^3$. Нормированный неотрицательный цвет определяется как
\begin{equation}
\tilde{\mathbf{c}}_i=
\operatorname{softplus}(\mathbf{z}_i^c),
\end{equation}
\begin{equation}
\mathbf{c}_i
=
\frac{\tilde{\mathbf{c}}_i}
{\|\tilde{\mathbf{c}}_i\|_1+\varepsilon}.
\end{equation}

Следовательно,
\begin{equation}
c_{i,R}+c_{i,G}+c_{i,B}\approx 1,
\qquad
c_{i,k}\geq 0.
\end{equation}

Интенсивность $\eta_i$ отвечает за общую энергию ячейки, а $\mathbf{c}_i$ --- за её хроматичность. Это устраняет масштабную неоднозначность между яркостью и цветом, возникающую при независимой оптимизации всех RGB-компонент вместе с интенсивностью:
\begin{equation}
\eta_i\mathbf{c}_i
\neq
(\alpha\eta_i)
\left(
\frac{\mathbf{c}_i}{\alpha}
\right)
\end{equation}
после введения нормировки цвета.

\subsection{Параметр угловой ширины}

Параметр $\theta_i$ ограничивается диапазоном
\begin{equation}
\theta_{\min}
\leq
\theta_i
\leq
\theta_{\max},
\qquad
0<\theta_{\min}<\theta_{\max}\leq\frac{\pi}{2}.
\end{equation}

Для этого вводится неограниченная переменная $z_i^\theta$:
\begin{equation}
\theta_i
=
\theta_{\min}
+
\left(\theta_{\max}-\theta_{\min}\right)
\sigma(z_i^\theta),
\end{equation}
где $\sigma(\cdot)$ --- сигмоида.

Параметр $\theta_i$ интерпретируется как эффективная угловая ширина локальной диаграммы излучения LED-элемента. Он не обязан в точности совпадать с паспортным углом половинной мощности отдельного светодиода: в случае матовой крышки он описывает эффективную направленность участка панели после рассеяния.

\section{Модель излучения и переноса света}

\subsection{Косинусная диаграмма направленности}

Для точки $\mathbf{x}$ сцены и ячейки источника $i$ определяется направление распространения света:
\begin{equation}
\boldsymbol{\omega}_{i\rightarrow x}
=
\frac{\mathbf{x}-\mathbf{x}_i}
{\|\mathbf{x}-\mathbf{x}_i\|}.
\end{equation}

Поскольку нормаль источника направлена вниз, направление $\boldsymbol{\omega}_{i\rightarrow x}$ для освещаемых поверхностей согласовано с $\mathbf{n}_L$. Угол отклонения от оптической оси источника равен
\begin{equation}
\alpha_i(\mathbf{x})
=
\arccos\left(
\operatorname{clamp}
\left(
\mathbf{n}_L\cdot\boldsymbol{\omega}_{i\rightarrow x},
-1,1
\right)
\right).
\end{equation}

Косинусная модель направленности имеет вид
\begin{equation}
D_i(\boldsymbol{\omega}_{i\rightarrow x})
=
\max\left(
0,\,
\cos\alpha_i(\mathbf{x})
\right)^{\kappa_i}.
\end{equation}

С учётом скалярного произведения это выражение можно записать как
\begin{equation}
D_i(\boldsymbol{\omega}_{i\rightarrow x})
=
\max\left(
0,\,
\mathbf{n}_L\cdot\boldsymbol{\omega}_{i\rightarrow x}
\right)^{\kappa_i}.
\end{equation}

Параметр $\kappa_i$ связывается с восстанавливаемым углом $\theta_i$. Если $\theta_i$ определяется как угол половинной мощности, то
\begin{equation}
D_i(\theta_i)=\frac{1}{2},
\end{equation}
откуда следует
\begin{equation}
\left(\cos\theta_i\right)^{\kappa_i}=\frac{1}{2},
\end{equation}
\begin{equation}
\kappa_i
=
\frac{\ln(1/2)}{\ln(\cos\theta_i)}.
\end{equation}

Таким образом,
\begin{equation}
D_i(\boldsymbol{\omega})
=
\max\left(
0,\mathbf{n}_L\cdot\boldsymbol{\omega}
\right)^{
\frac{\ln(1/2)}{\ln(\cos\theta_i)}
}.
\end{equation}

Данная модель гладкая по направлению и не содержит бинарной границы конуса. Поэтому она более удобна для дифференцируемой оптимизации, чем жёсткое отсечение по углу.

\subsection{Локальная радиантность}

Радиантность, испускаемая ячейкой источника $i$ в направлении $\boldsymbol{\omega}$, определяется как
\begin{equation}
\mathbf{L}_{e,i}(\boldsymbol{\omega})
=
\eta_i\,
\mathbf{c}_i\,
D_i(\boldsymbol{\omega}).
\end{equation}

Физически корректная формулировка может включать нормировочный коэффициент $Z(\kappa_i)$, сохраняющий полную мощность ячейки при изменении направленности:
\begin{equation}
\mathbf{L}_{e,i}(\boldsymbol{\omega})
=
\eta_i\,
\mathbf{c}_i\,
Z(\kappa_i)
\max\left(0,\mathbf{n}_L\cdot\boldsymbol{\omega}\right)^{\kappa_i}.
\end{equation}

Для осесимметричной диаграммы на полусфере коэффициент может быть выбран как
\begin{equation}
Z(\kappa_i)
=
\frac{\kappa_i+1}{2\pi},
\end{equation}
если $\eta_i$ интерпретируется как интегральная мощность на единицу площади либо как пропорциональная ей величина. Тогда выполняется нормировка
\begin{equation}
\int_{\Omega^+}
Z(\kappa_i)
\cos(\alpha)^{\kappa_i}\,d\omega
=1.
\end{equation}

При иной фотометрической интерпретации $\eta_i$ можно трактовать как пиковую радиантность ячейки и не использовать $Z(\kappa_i)$. В описании метода необходимо выбрать одну интерпретацию и не смешивать их: при наличии $Z(\kappa_i)$ интенсивность означает энергию, интегрированную по направлениям; без него она означает амплитуду диаграммы в направлении нормали. Для реконструкции LED-источников предпочтительно использовать нормированный вариант, поскольку он уменьшает неоднозначность между увеличением интенсивности $\eta_i$ и сужением диаграммы направленности.

\subsection{Материалы сцены}

Пол и стены моделируются как нейтральные ламбертовы поверхности с известным серым альбедо:
\begin{equation}
\boldsymbol{\rho}(\mathbf{x})
=
\rho
\begin{bmatrix}
1\\
1\\
1
\end{bmatrix},
\qquad
0<\rho<1.
\end{equation}

Их BRDF равна
\begin{equation}
f_r(\mathbf{x},\boldsymbol{\omega}_i,\boldsymbol{\omega}_o)
=
\frac{\rho}{\pi}.
\end{equation}

Предполагается отсутствие собственной эмиссии, зеркального отражения, пространственной вариации альбедо и выраженных субповерхностных эффектов. Нейтральный серый материал обеспечивает одинаковый отклик в RGB-каналах и тем самым позволяет отделять цвет света от цвета отражающей поверхности.

\subsection{Трассировка путей с двумя отскоками}

Наблюдаемое HDR-значение в пикселе рассчитывается методом трассировки путей с максимум двумя отражательными отскоками после попадания луча камеры в сцену.

Пусть луч камеры для пикселя $p$ впервые пересекает сцену в точке $\mathbf{x}_0$ с нормалью $\mathbf{n}_0$. Тогда учитываются следующие классы световых путей:
\begin{enumerate}
    \item Прямое освещение, без промежуточных отражений:
    \begin{equation}
    \text{camera}
    \rightarrow
    \mathbf{x}_0
    \rightarrow
    A_L.
    \end{equation}

    \item Одно межотражение:
    \begin{equation}
    \text{camera}
    \rightarrow
    \mathbf{x}_0
    \rightarrow
    \mathbf{x}_1
    \rightarrow
    A_L,
    \end{equation}
    где $\mathbf{x}_1$ лежит на полу или одной из стен.

    \item Два межотражения:
    \begin{equation}
    \text{camera}
    \rightarrow
    \mathbf{x}_0
    \rightarrow
    \mathbf{x}_1
    \rightarrow
    \mathbf{x}_2
    \rightarrow
    A_L.
    \end{equation}
\end{enumerate}

Следовательно, ограничение в два отскока относится к двум последовательным диффузным отражениям света в сцене. Непосредственное попадание из точки поверхности на источник не считается отражательным отскоком.

Для пути $\mathcal{P}$ длины $K$ со взаимодействиями $\mathbf{x}_0,\ldots,\mathbf{x}_{K-1}$ и завершающим попаданием на ячейку источника $i$ его вклад имеет общий вид
\begin{equation}
\mathbf{C}_{\mathcal{P}}
=
\mathbf{L}_{e,i}(\boldsymbol{\omega}_L)
\prod_{j=0}^{K-1}
\left[
\frac{\rho}{\pi}
\frac{
\left|
\mathbf{n}_j\cdot
\boldsymbol{\omega}_{j\rightarrow j+1}
\right|
}{
p_j(\boldsymbol{\omega}_{j\rightarrow j+1})
}
\right]
W_{\mathrm{geom}}(\mathcal{P}),
\end{equation}
где $p_j$ --- плотность вероятности семплирования направления на $j$-м отражении; $W_{\mathrm{geom}}$ включает геометрические члены, видимость и расстояния между вершинами пути; для ламбертовых поверхностей $f_r=\rho/\pi$.

В практической реализации для снижения дисперсии применяется next-event estimation: на каждом диффузном взаимодействии явно семплируется точка на источнике. Для дифференцируемой реконструкции целесообразно использовать общие случайные числа, то есть фиксированные последовательности случайных чисел между итерациями оптимизации. Это уменьшает изменение Монте--Карло шума между соседними состояниями параметров и стабилизирует оценку градиента.

\section{HDR-наблюдения и функция потерь}

\subsection{Представление HDR-данных}

Входные наблюдения представляют собой HDR-изображения
\begin{equation}
\mathbf{H}_m^\ast \in \mathbb{R}_+^{H\times W\times 3},
\qquad
m=1,\ldots,M,
\end{equation}
линейно пропорциональные радиантности, приходящей в камеру. Камерная обработка --- гамма-коррекция, тональное отображение, автоматическая экспозиция, автоматический баланс белого и клиппинг --- должна быть исключена либо устранена до оптимизации.

HDR-данные могут быть получены одним из двух способов:
\begin{itemize}
    \item из линейного RAW-кадра при отсутствии насыщения и достаточном отношении сигнал/шум;
    \item из серии LDR- или RAW-кадров с различными выдержками посредством радиометрической калибровки и HDR-слияния.
\end{itemize}

При многокадровом HDR-захвате короткие выдержки сохраняют информацию в ярких областях, а длинные --- в слабых отражениях на серых стенах и полу. Это позволяет покрыть динамический диапазон сцены без утраты данных из-за насыщения или недоэкспонирования.

\subsection{Предсказанное HDR-изображение}

Для каждой камеры $m$ дифференцируемый рендерер формирует линейное HDR-изображение:
\begin{equation}
\hat{\mathbf{H}}_m
=
s_m\,
\mathcal{R}
\left(
\mathcal{S},
\mathcal{G}_L,
\mathbf{T};
\mathcal{C}_m
\right),
\end{equation}
где $\mathcal{C}_m$ --- известные параметры камеры, а $s_m>0$ --- глобальный коэффициент радиометрического масштаба.

Если абсолютная радиометрическая калибровка камеры и источника выполнена, $s_m$ фиксируется. Если наблюдения известны только с точностью до масштаба, он параметризуется как
\begin{equation}
s_m=\exp(q_m)
\end{equation}
и оптимизируется совместно с параметрами источника. При одной камере и отсутствии абсолютной калибровки необходимо зафиксировать одну масштабную степень свободы, например
\begin{equation}
\sum_i\eta_i\Delta A_i=P_0,
\end{equation}
либо установить $s_1=1$.

\subsection{Фотометрическая HDR-функция потерь}

Для предотвращения доминирования ярких областей прямого освещения и сохранения чувствительности к слабым межотражениям используются два комплементарных члена.

Ошибка в линейном пространстве определяется как
\begin{equation}
\mathcal{L}_{\mathrm{lin}}
=
\frac{1}{M}
\sum_{m=1}^{M}
\frac{
\sum_{p\in\Omega_m}
M_{m,p}
\left\|
\hat{\mathbf{H}}_{m,p}
-
\mathbf{H}_{m,p}^\ast
\right\|_1
}{
\sum_{p\in\Omega_m}M_{m,p}+\varepsilon
}.
\end{equation}

Ошибка в логарифмическом HDR-пространстве определяется как
\begin{equation}
\mathcal{L}_{\log}
=
\frac{1}{M}
\sum_{m=1}^{M}
\frac{
\sum_{p\in\Omega_m}
M_{m,p}
\left\|
\log\left(\hat{\mathbf{H}}_{m,p}+\varepsilon\right)
-
\log\left(\mathbf{H}_{m,p}^\ast+\varepsilon\right)
\right\|_1
}{
\sum_{p\in\Omega_m}M_{m,p}+\varepsilon
}.
\end{equation}

Маска достоверности $M_{m,p}\in[0,1]$ исключает или ослабляет вклад пикселей, насыщенных во всех кадрах исходного брекетинга, находящихся ниже уровня достоверного сигнала, не принадлежащих геометрии пола или стен, искажённых flare, ghosting или ошибками совмещения HDR-кадров, а также расположенных в областях с недостаточно известной геометрией или материалом.

Логарифмическая составляющая необходима вследствие широкого диапазона яркостей HDR-представления: она не позволяет нескольким ярким пикселям полностью подавить градиент от слабых, но информативных межотражений.

\subsection{Регуляризация текстурных полей}

Поскольку число свободных параметров источника может существенно превышать число независимо наблюдаемых световых эффектов, задача нуждается в пространственных априорных ограничениях.

Для интенсивности используется TV-регуляризация:
\begin{equation}
\mathcal{L}_{\mathrm{TV}}^\eta
=
\sum_i
\left(
|\eta_{i+\mathbf{e}_u}-\eta_i|
+
|\eta_{i+\mathbf{e}_v}-\eta_i|
\right).
\end{equation}

Для параметра направленности:
\begin{equation}
\mathcal{L}_{\mathrm{TV}}^\theta
=
\sum_i
\left(
|\theta_{i+\mathbf{e}_u}-\theta_i|
+
|\theta_{i+\mathbf{e}_v}-\theta_i|
\right).
\end{equation}

Для цветовой текстуры:
\begin{equation}
\mathcal{L}_{\mathrm{TV}}^c
=
\sum_i
\left(
\|\mathbf{c}_{i+\mathbf{e}_u}-\mathbf{c}_i\|_1
+
\|\mathbf{c}_{i+\mathbf{e}_v}-\mathbf{c}_i\|_1
\right).
\end{equation}

Их взвешенная сумма имеет вид
\begin{equation}
\mathcal{L}_{\mathrm{smooth}}
=
\lambda_\eta\mathcal{L}_{\mathrm{TV}}^\eta
+
\lambda_\theta\mathcal{L}_{\mathrm{TV}}^\theta
+
\lambda_c\mathcal{L}_{\mathrm{TV}}^c.
\end{equation}

Эта регуляризация имеет физическую интерпретацию. Для светодиодной панели с матовой крышкой соседние области эффективной поверхности обычно имеют близкие характеристики излучения. Для открытой LED-матрицы регуляризация может быть слабее, чтобы не стирать различия между отдельными светодиодными модулями.

\subsection{Ограничение полной мощности}

Если $\eta_i$ интерпретируется как интегральная мощность локальной ячейки, может использоваться энергетическое ограничение:
\begin{equation}
\mathcal{L}_{\mathrm{power}}
=
\left(
\sum_i\eta_i\Delta A_i-P_0
\right)^2.
\end{equation}

Оно полезно при известной номинальной мощности источника или для устранения глобальной масштабной неоднозначности между интенсивностью света и радиометрическим масштабом камеры.

\subsection{Итоговая функция потерь}

Полная функция оптимизации задаётся как
\begin{equation}
\mathcal{L}
=
\lambda_{\mathrm{lin}}\mathcal{L}_{\mathrm{lin}}
+
\lambda_{\log}\mathcal{L}_{\log}
+
\lambda_{\mathrm{smooth}}\mathcal{L}_{\mathrm{smooth}}
+
\lambda_{\mathrm{power}}\mathcal{L}_{\mathrm{power}}.
\end{equation}

Искомые параметры определяются выражением
\begin{equation}
\hat{\boldsymbol{\pi}}
=
\arg\min_{\boldsymbol{\pi}}
\mathcal{L}
\left(
\hat{\mathbf{H}}(\boldsymbol{\pi}),
\mathbf{H}^\ast
\right).
\end{equation}

\section{Coarse-to-fine реконструкция}

Прямая реконструкция в высоком разрешении плохо обусловлена: высокочастотные вариации текстуры источника могут приводить к практически неразличимым изменениям в наблюдаемом освещении. Поэтому применяется многоуровневая coarse-to-fine оптимизация.

Пусть последовательность разрешений источника задана как
\begin{equation}
\mathcal{Q}=
\left\{
8\times 8,\;
16\times 16,\;
32\times 32,\;
64\times 64
\right\}.
\end{equation}

Начальная реконструкция проводится для низкого разрешения $8\times 8$, где одна ячейка соответствует относительно крупной области LED-панели. Затем после сходимости найденные неограниченные параметры
\begin{equation}
z^\eta,\qquad z^\theta,\qquad \mathbf{z}^c
\end{equation}
билинейно интерполируются к следующему разрешению. Оптимизация продолжается с этой инициализацией.

Важно интерполировать именно неограниченные оптимизационные переменные, а не уже нормированные физические поля. В частности,
\begin{equation}
z^\eta_{\ell+1}
=
\operatorname{Upsample}(z^\eta_\ell),
\end{equation}
\begin{equation}
z^\theta_{\ell+1}
=
\operatorname{Upsample}(z^\theta_\ell),
\end{equation}
\begin{equation}
\mathbf{z}^{c}_{\ell+1}
=
\operatorname{Upsample}(\mathbf{z}^{c}_\ell).
\end{equation}

После интерполяции снова применяются softplus, сигмоида и цветовая нормировка. Это предотвращает нарушение ограничений
\begin{equation}
\eta_i\geq 0,
\qquad
\theta_i\in[\theta_{\min},\theta_{\max}],
\qquad
\|\mathbf{c}_i\|_1=1.
\end{equation}

На низких разрешениях могут использоваться более сильные коэффициенты регуляризации гладкости. По мере повышения разрешения коэффициенты $\lambda_\eta$, $\lambda_\theta$ и $\lambda_c$ постепенно уменьшаются, позволяя восстанавливать детали, поддерживаемые наблюдениями.

\section{Алгоритм оптимизации}

В качестве оптимизатора используется Adam или AdamW. Один шаг на уровне разрешения $\ell$ состоит из следующих операций:
\begin{enumerate}
    \item Получение физических параметров источника:
    \begin{equation}
    \eta=\operatorname{softplus}(z^\eta),
    \end{equation}
    \begin{equation}
    \theta=
    \theta_{\min}
    +(\theta_{\max}-\theta_{\min})\sigma(z^\theta),
    \end{equation}
    \begin{equation}
    \mathbf{c}=
    \frac{\operatorname{softplus}(\mathbf{z}^c)}
    {\|\operatorname{softplus}(\mathbf{z}^c)\|_1+\varepsilon}.
    \end{equation}

    \item Преобразование $\theta_i$ в показатель косинусной диаграммы:
    \begin{equation}
    \kappa_i
    =
    \frac{\ln(1/2)}{\ln(\cos\theta_i)}.
    \end{equation}

    \item Выполнение дифференцируемой трассировки путей с двумя межотражениями и получение $\hat{\mathbf{H}}_m$ для всех камер.

    \item Вычисление фотометрических HDR-функций потерь и регуляризационных членов.

    \item Вычисление градиентов:
    \begin{equation}
    \nabla_{z^\eta,z^\theta,\mathbf{z}^c}\mathcal{L}.
    \end{equation}

    \item Обновление параметров оптимизатором.

    \item После сходимости на текущем уровне --- увеличение разрешения текстуры и продолжение оптимизации.
\end{enumerate}

% Псевдокод алгоритма реконструкции:
%
% Input:
%     HDR-наблюдения H*
%     известная геометрия сцены S
%     известная геометрия источника GL
%     параметры камер C
%     последовательность разрешений Q
%
% Initialize:
%     zη ← constant intensity field
%     zθ ← constant angular-width field
%     zc ← neutral-white color field
%
% for q in Q:
%     zη, zθ, zc ← UpsampleToResolution(zη, zθ, zc, q)
%
%     for t = 1 ... Kq:
%         η ← softplus(zη)
%         θ ← θmin + (θmax - θmin) sigmoid(zθ)
%         κ ← log(1/2) / log(cos(θ))
%         c ← softplus(zc) / (sum_channels(softplus(zc)) + ε)
%
%         H_hat ← PathTrace(
%             scene=S,
%             light_geometry=GL,
%             intensity=η,
%             color=c,
%             directional_exponent=κ,
%             max_diffuse_bounces=2
%         )
%
%         L ← λlin Llin(H_hat, H*)
%            + λlog Llog(H_hat, H*)
%            + λsmooth Lsmooth(η, θ, c)
%            + λpower Lpower(η)
%
%         zη, zθ, zc ← AdamStep(zη, zθ, zc, ∇L)
%
% Output:
%     T_hat = [θ, η, cR, cG, cB]

\section{Практические замечания и ограничения}

\subsection{Интерпретация параметров LED-источника}

Восстанавливаемая текстура не обязательно совпадает с дискретной физической картой отдельных светодиодов. Она представляет эффективную модель источника в масштабе наблюдаемости.

Для открытой LED-матрицы локальные максимумы $\eta(u,v)$ могут соответствовать отдельным LED-модулям или группам диодов. Для панели с матовым диффузором пространственная эмиссия сглаживается, поэтому реконструируемая карта описывает распределение света на внешней поверхности рассеивателя.

При расстоянии между источником и сценой, значительно превышающем шаг светодиодной решётки, влияние отдельных LED на освещённость дополнительно интегрируется геометрией переноса света. В этом режиме восстанавливаемы преимущественно низкочастотные вариации интенсивности, цвета и направленности.

\subsection{Влияние двух межотражений}

Учёт двух диффузных отражений особенно важен для угловой сцены. Свет может попасть на пол, отразиться на стену, затем --- на вторую стену или обратно на пол. Например, учитываются пути
\begin{equation}
A_L
\rightarrow
S_f
\rightarrow
S_x
\rightarrow
\text{camera}
\end{equation}
и
\begin{equation}
A_L
\rightarrow
S_x
\rightarrow
S_y
\rightarrow
\text{camera}.
\end{equation}

Такие пути вносят вклад в тёмные области и повышают информативность наблюдений относительно распределения света на источнике.

При малом альбедо $\rho$ вклад второго межотражения ограничен порядком $\rho^2$, однако для HDR-наблюдений он может оставаться измеримым. Включение этих путей уменьшает систематическую ошибку, которую дала бы модель только прямого освещения.

\subsection{Неоднозначности}

Даже при фиксированной геометрии остаются несколько фундаментальных неоднозначностей:
\begin{itemize}
    \item При неизвестном абсолютном масштабе HDR-радиантности невозможно независимо определить глобальный масштаб интенсивности $\eta$ и множитель $s_m$.
    \item Узкая яркая диаграмма и более широкая, но слабая диаграмма могут порождать похожую освещённость в ограниченной области наблюдения.
    \item Высокочастотная структура LED-матрицы может быть неразличима из-за расстояния до сцены, интегрирования по площади, ограниченного разрешения камеры и межотражений.
    \item Ошибки геометрии источника, положения камеры, альбедо или HDR-калибровки могут компенсироваться ложными вариациями восстановленной текстуры.
\end{itemize}

Поэтому получаемый результат следует интерпретировать как регуляризованную оценку пространственно изменяющегося излучения, наблюдаемо эквивалентную реальному источнику для заданной сцены, камер и режима рендеринга.





\endinput